% !TeX program = lualatex
% This document requires LuaLaTeX for compilation

\documentclass[a4paper,12pt]{article}

% Engine detection and compatibility check
\RequirePackage{iftex}
\ifLuaTeX
  % Document is being compiled with LuaTeX - proceed normally
\else
  \PackageError{main}{This document requires LuaLaTeX for compilation.
    Please use: lualatex main.tex}
    {This document uses fontspec and other LuaTeX-specific features.}
\fi

\usepackage{fontspec}
\usepackage[english]{babel}

% Font configuration for LuaTeX
\setmainfont{Latin Modern Roman}
\setsansfont{Latin Modern Sans}
\setmonofont{Latin Modern Mono}

\usepackage{amsmath} % Para equações matemáticas
\usepackage{geometry}
\geometry{margin=1in}
\usepackage{hyperref}
\usepackage[sort,numbers]{natbib}
\usepackage{parskip}
\usepackage{longtable}
\usepackage{xcolor}
\usepackage{tabularx}
\setlength{\marginparwidth}{2cm}
\usepackage[colorinlistoftodos,prependcaption,textsize=tiny]{todonotes}
\usepackage{comment}
\usepackage{graphicx} % Para resizebox
\usepackage{array} % Para colunas com largura fixa
\usepackage{listings}
\usepackage{siunitx} % For \si{\micro\second}
\usepackage[acronym,nonumberlist]{glossaries}
\usepackage{placeins}

% Configuring listings for code and log display
\lstset{
  basicstyle=\ttfamily\small,
  breaklines=true,
  breakatwhitespace=true,
  frame=single,
  numbers=left,
  numberstyle=\tiny,
  keywordstyle=\color{blue},
  commentstyle=\color{gray},
  showstringspaces=false,
  tabsize=2,
  extendedchars=true
}


\title{Lista 1 - Modelagem de Conhecimento em Lógica Aplicada à Análise de Cafés Especiais}
\author{André Thiago Pfleger \\  Gustavo Girotto \\ João Pedro Schmidt Cordeiro}
\date{\today}

\begin{document}

\begin{titlepage}
   \begin{center}
      \textsc{Universidade Federal de Santa Catarina} \\
      \textsc{Centro Tecnológico} \\
      \textsc{Departamento de Informática e Estatística} \\[4cm]

      \textbf{\Large{INE5430 - Inteligência Artificial}} \\[2cm]

      \textbf{\huge{Trabalho 1 - Construção de um Sistema Especialista para Análise de Cafés Especiais}} \\[3cm]

      André Thiago Pfleger \\
      Gustavo Girotto \\
      João Pedro Schmidt Cordeiro \\
      
      \vfill
      
      {Florianópolis} \\
      {\today}
   \end{center}
\end{titlepage}

\section{Introdução}

Este trabalho apresenta o desenvolvimento de um sistema especialista em CLIPS para análise e recomendação de harmonizações gastronômicas com cafés especiais. O sistema utiliza uma base de conhecimento estruturada contendo 12 variedades distintas de café, cada uma com características específicas de torra, moagem e métodos de preparo recomendados.

O objetivo principal é demonstrar a aplicação prática de sistemas baseados em conhecimento para domínios especializados, implementando regras de inferência que correlacionam características técnicas do café com harmonizações gastronômicas apropriadas. O sistema foi desenvolvido seguindo princípios de engenharia de conhecimento, com ênfase na consistência, determinismo e cobertura completa do domínio.

A abordagem utilizada combina frames (via \texttt{deftemplate} do CLIPS) para representação estruturada do conhecimento com regras de produção para inferência automática de recomendações. Durante o desenvolvimento, foram implementados mecanismos de controle para garantir que cada variedade de café receba exatamente uma harmonização de cada categoria (pão, queijo e bolo), resolvendo conflitos através de sistemas de prioridade bem definidos.



\section{Variedades de café (descrição, torra, moagem e preparo)}
A seguir, para cada variedade há: descrição resumida, torra recomendada (e justificativa), moagem sugerida e modos de preparo indicados.

\subsection{Robusta (Coffea canephora)}
{Descrição} Sabor acentuado e rústico, corpo pesado, maior resistência a doenças e teor de cafeína elevado; frequentemente utilizado em blends e produtos solúveis.  
{Torra recomendada} \textbf{Média a escura} — realça corpo e notas torradas, equilibrando a acidez naturalmente baixa.  
{Moagem sugerida} Fina (espresso) a média (filtros rápidos), conforme o método.  
{Modos de preparo indicados} Espresso, moka, blends industriais e cafés instantâneos.

\subsection{Conilon (grupo canephora, termo comum no Brasil)}
{Descrição} Variante de canephora muito explorada no Brasil; grãos tipicamente menores, amargor pronunciado e alto teor de cafeína; uso comum em commodity e blends.  
{Torra recomendada} \textbf{Média a escura} — enfatiza corpo e notas de torra.  
{Moagem sugerida} Fina (espresso) / média (filtrados comerciais).  
{Modos de preparo indicados} Espresso, blends e solúveis.

\subsection{Arábica (Coffea arabica — genérico)}
{Descrição} A variedade mais consumida globalmente; perfis sensoriais complexos com acidez, doçura e notas frutadas/ florais, fortemente dependentes de terroir e altitude.  
{Torra recomendada} \textbf{Clara a média} — preserva as características de origem e a acidez desejada.  
{Moagem sugerida} Média a média-fina (pour-over); fina para espresso specialty.  
{Modos de preparo indicados} Pour-over (V60), Chemex, Aeropress, filtros e espresso specialty.

\subsection{Bourbon}
{Descrição} Subgrupo de Arábica com doçura pronunciada e aroma intenso; bebida tipicamente suave.  
{Torra recomendada} \textbf{Clara a média} — evidencia doçura e nuances carameladas/frutadas.  
{Moagem sugerida} Média a média-fina.  
{Modos de preparo indicados} Pour-over, Chemex, Aeropress; prensa francesa quando se busca maior corpo.

\subsection{Catuaí (Catuai)}
{Descrição} Cultivar brasileira resultante de cruzamentos; planta compacta, produtividade e doçura natural; amplamente plantado.  
{Torra recomendada} \textbf{Média} — equilíbrio entre doçura e corpo.  
{Moagem sugerida} Média.  
{Modos de preparo indicados} Filtrados (V60, coador), blends para espresso, prensa francesa.

\subsection{Mundo Novo (Novo Mundo)}
{Descrição} Cruzamento natural (Typica × Bourbon); vigoroso, produtivo; tende a oferecer corpo mais pronunciado e doçura.  
{Torra recomendada} \textbf{Média a média-escura} — realça corpo e notas doces.  
{Moagem sugerida} Média a média-grosseira.  
{Modos de preparo indicados} Pour-over, prensa francesa, blends que visem corpo.

\subsection{Caturra}
{Descrição} Mutação do Bourbon com planta de porte compacto; boa produtividade inicial; produz cafés de boa qualidade e doçura.  
{Torra recomendada} \textbf{Clara a média} — preserva equilíbrio entre acidez e doçura.  
{Moagem sugerida} Média.  
{Modos de preparo indicados} Filtrados, Chemex, Aeropress, single-origin.

\subsection{Acaiá (Acaia)}
{Descrição} Mutação do Mundo Novo; variedade rara, com bom desempenho em altitudes acima de \(\sim\)800\,m; xícara com notas frutadas e lembranças de chocolate.  
{Torra recomendada} \textbf{Clara a média} — evidencia notas frutadas e acidez média.  
{Moagem sugerida} Média-fina (pour-over) / média para métodos com mais contato.  
{Modos de preparo indicados} Pour-over, Aeropress, cupping.

\subsection{Typica}
{Descrição} Variedade histórica e clássica; perfil doce e floral/complexo; menor produtividade e maior sensibilidade a pragas.  
{Torra recomendada} \textbf{Clara} — preserva delicadeza e complexidade aromática.  
{Moagem sugerida} Média-fina.  
{Modos de preparo indicados} Pour-over, V60, cupping.

\subsection{Geisha (Gesha)}
{Descrição} Variedade muito expressiva aromaticamente (jasmim, flores, frutas de caroço); extremamente valorizada no mercado specialty.  
{Torra recomendada} \textbf{Clara} (por vezes muito clara) — para não mascarar notas voláteis e florais.  
{Moagem sugerida} Média-fina a média.  
{Modos de preparo indicados} Pour-over (V60), Kalita/Wave, métodos que priorizem aroma e limpeza.

\subsection{Maragogipe}
{Descrição} Mutação de Typica com sementes excepcionalmente grandes ("grão-elefante"); perfil muitas vezes suave e frutado; baixa produtividade.  
{Torra recomendada} \textbf{Clara a média} — destaca notas delicadas.  
{Moagem sugerida} Média.  
{Modos de preparo indicados} Filtrados especiais, single-origin.

\subsection{Pacamara (Pacas × Maragogipe)}
{Descrição} Híbrido que produz grãos grandes e perfil complexo (floral, frutado e corpo médio).  
{Torra recomendada} \textbf{Clara a média} — realça complexidade sem sobrecarregar com notas de torra.  
{Moagem sugerida} Média-fina a média.  
{Modos de preparo indicados} Pour-over, Aeropress, extrações experimentais em pequenos lotes.

\section{Recomendações práticas gerais}
\begin{itemize}
  \item \textbf{Moagem por método (resumo):} Espresso: \emph{fina}; Aeropress/V60: \emph{média-fina}; Chemex: \emph{média-grosseira}; French press: \emph{grossa}; Cold brew: \emph{extra grossa}.
  \item \textbf{Torra:} Clara preserva origem e acidez; média equilibra doçura e corpo; escura traz notas de torra e reduz percepção de acidez — escolha conforme objetivo sensorial.
  \item \textbf{Ajustes (dial-in):} Afine moagem, dose e tempo/temperatura antes de alterar torra; pequenas mudanças na moagem provocam grandes mudanças na extração.
\end{itemize}

\section{Variedades de comidas selecionadas}

\subsection{Pães Integrais}
Descrição: Pães feitos com farinha integral, sem adição de fermento químico.
\begin{itemize}
  \item \textbf{Com sementes} — Textura mais densa e crocante, com notas terrosas que complementam cafés de corpo mais acentuado e torras mais escuras. As sementes (girassol, abóbora, linhaça) adicionam óleos naturais que equilibram a intensidade do café.
  \item \textbf{Fermentação Natural} — Sabor mais complexo e levemente ácido devido ao processo de fermentação longa, criando harmonia com cafés de torras claras que preservam a acidez natural do grão. A textura arejada contrasta com a delicadeza dos aromas florais.
  \item \textbf{Com castanhas e Frutas secas} — Doçura natural e notas caramelizadas das frutas secas (damasco, tâmara) e oleosidade das castanhas (nozes, amêndoas) que complementam cafés de torra média, realçando suas características equilibradas entre acidez e corpo.
\end{itemize}

\subsection{Queijos semicurados}
Descrição: Queijos feitos com leite de vaca, sem adição de conservantes.
\begin{itemize}
  \item \textbf{Parmesão Jovem} — Sabor mais suave que o parmesão envelhecido, com notas lácteas e levemente salgadas. Sua textura firme mas não ressecada harmoniza com cafés de moagem fina, onde a extração mais intensa encontra equilíbrio na cremosidade residual do queijo.
  \item \textbf{Gouda} — Queijo de sabor adocicado e textura cremosa, com notas de caramelo que se desenvolvem durante a cura. Complementa perfeitamente cafés de moagem média-fina, onde a doçura natural do queijo realça as notas frutadas e equilibra possíveis amargor residual.
  \item \textbf{Canastra Meia Cura} — Queijo artesanal brasileiro com sabor mais acentuado e notas picantes características. Sua personalidade marcante harmoniza com cafés de moagem média, criando um contraste interessante entre a rusticidade do queijo e a suavidade da extração.
\end{itemize}

\subsection{Bolos Cítricos}
Descrição: Bolos feitos com farinha de trigo, sem adição de fermento químico.
\begin{itemize}
  \item \textbf{Laranja com calda} — Bolo úmido com acidez equilibrada da laranja e doçura da calda, criando complexidade de sabores. A calda adiciona textura e intensifica os aromas cítricos, harmonizando perfeitamente com métodos de filtragem que preservam as notas delicadas e aromáticas do café.
  \item \textbf{Limão} — Acidez mais pronunciada e refrescante, com notas cítricas intensas que despertam o paladar. Sua característica marcante complementa métodos de pressão (espresso), onde a intensidade do café encontra contraponto na vivacidade ácida do limão, criando equilíbrio gustativo.
  \item \textbf{Tangerina com cobertura leve} — Sabor cítrico mais suave e doce, com cobertura que adiciona cremosidade sem mascarar os aromas. Harmoniza idealmente com métodos de imersão, onde o tempo de contato prolongado permite que as notas sutis do café se integrem com a delicadeza da tangerina.
\end{itemize}


% construindo a tabela com base nas descrições
\section{Sistema Especialista: Desenvolvimento e Refinamentos}

O sistema especialista desenvolvido em CLIPS utiliza uma abordagem baseada em frames (via \texttt{deftemplate}) combinada com regras de inferência para gerar recomendações personalizadas de harmonização para cada variedade de café.

\subsection{Arquitetura do Sistema}
O sistema é estruturado em três componentes principais:
\begin{itemize}
  \item \textbf{Templates (Frames):} Definem a estrutura dos fatos sobre variedades e recomendações
  \item \textbf{Fatos Iniciais:} Base de conhecimento com as 12 variedades de café estudadas
  \item \textbf{Regras de Inferência:} Conjunto de regras que geram recomendações baseadas nas características de cada variedade
\end{itemize}

\subsection{Lógica de Harmonização e Critérios de Desempate}
Durante o desenvolvimento, identificou-se a necessidade de implementar critérios rigorosos para garantir harmonizações consistentes. O sistema foi refinado para produzir exatamente \textbf{uma harmonização de cada categoria} (pão, queijo e bolo) por variedade de café.

\subsubsection{Mapeamento por Características}
As harmonizações são determinadas através de três mapeamentos independentes:
\begin{itemize}
  \item \textbf{Pães:} Baseados exclusivamente no \emph{tipo de torra} (mapeamento 1:1)
    \begin{itemize}
      \item Clara → Pães integrais de fermentação natural
      \item Clara a Média → Pães integrais de fermentação natural
      \item Média → Pães integrais com castanhas e frutas secas  
      \item Média a Média Escura → Pães integrais com sementes
      \item Média a Escura → Pães integrais com sementes
    \end{itemize}
  \item \textbf{Queijos:} Baseados exclusivamente no \emph{tipo de moagem}
    \begin{itemize}
      \item Fina a média → Queijo semicurado parmesão jovem
      \item Média a fina → Queijo semicurado gouda
      \item Média → Queijo semicurado canastra meia cura
      \item Média a Média-Grosseira → Queijo semicurado gouda
    \end{itemize}
  \item \textbf{Bolos:} Baseados no \emph{método de preparo} com sistema de prioridade
\end{itemize}

\subsubsection{Sistema de Prioridade para Desempate}
Para resolver conflitos quando uma variedade possui múltiplos métodos de preparo, implementou-se um sistema de prioridade usando \texttt{salience} no CLIPS:

\begin{enumerate}
  \item \textbf{Pressão} (prioridade 30) → Bolo cítrico de limão
  \item \textbf{Filtragem} (prioridade 20) → Bolo cítrico de laranja com calda  
  \item \textbf{Imersão} (prioridade 10) → Bolo cítrico de tangerina com cobertura leve
  \item \textbf{Comercial} (prioridade 5) → Bolo cítrico de limão
\end{enumerate}

Este sistema garante que, mesmo quando uma variedade suporta múltiplos métodos (como Mundo Novo: Imersão, Filtragem, Comercial), apenas um bolo seja selecionado com base na prioridade mais alta.

\subsection{Prevenção de Duplicações}
Cada regra de harmonização verifica se já existe qualquer item da mesma categoria antes de adicionar um novo, utilizando testes compostos com operadores lógicos:
\begin{lstlisting}[language=Lisp]
(test (not (or (member$ "Pão tipo 1" $?h)
               (member$ "Pão tipo 2" $?h)
               (member$ "Pão tipo 3" $?h))))
\end{lstlisting}

\subsection{Validação e Resultados}
O sistema foi testado e validado através da execução completa sobre todas as 12 variedades de café. O mapeamento segue rigorosamente a tabela de características, com cada regra correspondendo exatamente a uma linha específica da tabelade referência.

Os resultados demonstram que:
\begin{itemize}
  \item \textbf{Consistência:} Todas as variedades produzem exatamente 3 harmonizações (1 pão + 1 queijo + 1 bolo)
  \item \textbf{Determinismo:} O sistema de prioridade elimina ambiguidades, gerando sempre o mesmo resultado para a mesma entrada
  \item \textbf{Cobertura Completa:} Todas as 12 variedades recebem recomendações apropriadas baseadas em suas características específicas
  \item \textbf{Rastreabilidade:} Cada harmonização pode ser rastreada diretamente para uma linha específica da tabela
\end{itemize}

\subsubsection{Exemplos de Harmonizações Geradas}
O sistema produz recomendações como:
\begin{itemize}
  \item \textbf{Robusta} (Torra: Média a Escura, Moagem: Fina a média, Métodos: Pressão/Comercial) → Pães integrais com sementes + Queijo semicurado parmesão jovem + Bolo cítrico de limão
  \item \textbf{Geisha} (Torra: Clara, Moagem: Média a fina, Métodos: Filtragem) → Pães integrais de fermentação natural + Queijo semicurado gouda + Bolo cítrico de laranja com calda
  \item \textbf{Mundo Novo} (Torra: Média a Média Escura, Moagem: Média a Média-Grosseira, Métodos: Imersão/Filtragem/Comercial) → Pães integrais com sementes + Queijo semicurado gouda + Bolo cítrico de laranja com calda (Filtragem tem prioridade sobre Imersão e Comercial)
\end{itemize}

\section{Resumo das Variedades em Tabela}

% TODO: Aqui deve ser inserida a tabela com as variedades de café e as comidas selecionadas
\begin{table}[ht]
\centering
\begin{tabular}{l l l l}
\hline
\textbf{Variedade} & \textbf{Bolo} & \textbf{Pão} & \textbf{Queijo} \\
\hline
Robusta     & Com castanhas e frutas secas & Pães integrais                  & Queijos semicurados \\
Conilon     & Com castanhas e frutas secas & Pães integrais                  & Parmesão jovem      \\
Arábica     & Bolos cítricos               & Fermentação natural             & Queijos semicurados \\
Bourbon     & Com castanhas e frutas secas & Pães integrais com sementes     & Canastra            \\
Catuaí      & Bolos cítricos               & Pães integrais com sementes     & Meia Cura           \\
Mundo Novo  & Com castanhas e frutas secas & Pães integrais                  & Meia Cura           \\
Caturra     & Com castanhas e frutas secas & Fermentação natural             & Gouda               \\
Acaiá       & Bolos cítricos               & Pães integrais                  & Queijos semicurados \\
Typica      & Laranja com calda            & Fermentação natural             & Queijos semicurados \\
Geisha      & Limão                       & Fermentação natural             & Meia Cura           \\
Maragogipe  & Com castanhas e frutas secas & Pães integrais                  & Gouda               \\
Pacamara    & Tangerina com cobertura leve & Pães integrais com sementes     & Meia Cura           \\
\hline
\end{tabular}
\caption{Sugestões de harmonização (bolo, pão, queijo) por variedade}
\label{tab:harmonizacoes}
\end{table}

\FloatBarrier
\appendix
\section{Código Fonte em CLIPS}
\label{sec:codigo_clips}
O código a seguir representa a base de conhecimento completa desenvolvida em CLIPS para este trabalho.
\lstinputlisting[language=Lisp, caption={Base de Conhecimento em CLIPS sobre Variedades de Café}, label={lst:clips}]{main.clp}


\bibliographystyle{IEEEtran}
\bibliography{ref}
\nocite{counterculture_pourover}
\nocite{embrapa_catalog}
\nocite{perfectdailygrind_geisha}
\nocite{sca_brewing_chart}
\nocite{sca_standard_pdf}
\nocite{wcr_varieties}
\nocite{wikipedia_geisha}

\end{document}
